\chapter{Discussion} % Chapter 6
\label{Discussion}

% =============================================================================
% DISCUSSION (5-8 pages)
% Rubric: Clarity of Conclusions (10% of Report)
% =============================================================================

% -----------------------------------------------------------------------------
% SECTION 7.1: Interpretation of Results
% -----------------------------------------------------------------------------
\section{Interpretation of Results}

% TODO: 7.1.1 — Reasoning depth effects (IV1)
\subsection{Reasoning Depth and the Hobbesian Trap}
%   - Non-monotonic pattern: L0→L1→L2→L3 (if confirmed by Qwen results)
%   - Compare to Kuusela & Roy (AAMAS 2024):
%     Their finding: higher reasoning → more conflict (L0→L1→L2)
%     Our finding: [to be filled based on results]
%   - Key comparison: same fear spiral, different behavioral expression?
%     Kuusela & Roy: fear → pre-emptive attacks (only option in 2-action game)
%     Our game: fear → defensive hoarding (arm_self) or attacks, depending on params
%   - The fear spiral from deeper reasoning may be CONSERVED across game designs,
%     but behavioral expression depends on available actions and game structure

% TODO: 7.1.2 — Information structure effects (IV2)
\subsection{Information and Memory Effects}
%   - Does persistent memory change emergent structure?
%   - Hobbes's "diffidence": does local (incomplete) information increase arming?
%   - Do agents with memory develop more stable cooperative relationships?
%   - Connection to Harsanyi: incomplete information games in LLM agents

% TODO: 7.1.3 — Communication & enforcement effects (IV3)
\subsection{Communication and Enforcement}
%   - Does cheap talk actually change behavior? (Crawford & Sobel prediction)
%   - Does enforcement reduce conflict? (Hobbes's Leviathan hypothesis)
%   - Is enforcement MORE effective at higher reasoning depth?
%     (sophisticated agents may better exploit commitment mechanisms)

% TODO: 7.1.4 — Interaction effects
\subsection{Interaction Effects Across IVs}
%   - Does reasoning depth × information interact?
%     (does information matter MORE at higher reasoning depth?)
%   - Does reasoning depth × enforcement interact?
%     (does enforcement help MORE for deeper reasoners?)
%   - These interaction effects are the novel contribution

% -----------------------------------------------------------------------------
% SECTION 7.2: The Origins of Order
% -----------------------------------------------------------------------------
\section{The Origins of Order}
% TODO: Answer the title question
%   - Order emerges from the INTERACTION of reasoning, information, and enforcement
%   - Neither alone determines outcomes:
%     * Same structure, different reasoning → different outcomes
%     * Same reasoning, different information → different outcomes
%     * Same everything, with/without enforcement → different outcomes
%   - The thesis demonstrates that ALL THREE mechanism design inputs matter
%
% TODO: Phase transition interpretation
%   - If confirmed: critical point shifts with reasoning depth / information
%   - This is a genuine interaction between cognition/knowledge and structure
%   - Strongest theoretical contribution: parametric map of conditions under
%     which social order emerges vs collapses

% TODO: Connect to computational depth
%   \cite{pfau2024dots}: the manipulation works because it changes computation,
%     not just content. Even if traces are unfaithful, the prompt changes what
%     the model computes before acting.

% -----------------------------------------------------------------------------
% SECTION 7.3: On Reasoning Traces
% -----------------------------------------------------------------------------
\section{On Reasoning Traces}
% TODO: What can we claim?
%   - Traces correlate with behavior (L2 mentions opponent modeling, L0 doesn't)
%   - But correlation ≠ causation in trace content
%   - The PROMPT manipulation is causal (it's the experimental variable)
%   - The trace CONTENT may or may not reflect the actual computation
%   - Validation chapter results: what did truncation/anchors show?
%
% TODO: The "computation, not content" framing
%   - Our L0→L3 prompts change computational depth
%   - This is a real intervention even if traces are unfaithful
%   - Frame: "reasoning depth changes behavior" not "reasoning content explains behavior"
%
% TODO: Design choice — instruction, not computation budget
%   - We manipulate instructed reasoning depth, NOT computational budget
%   - Token counts are OBSERVED, not controlled
%   - If L2 and L3 produce similar token counts but different behavior,
%     the instruction matters beyond token count

% -----------------------------------------------------------------------------
% SECTION 7.4: Limitations
% -----------------------------------------------------------------------------
\section{Limitations}
% TODO: Model specificity
%   - Results on Qwen 3.5-27B — may not generalize to other architectures
%   - Gemma 2 results provide some cross-model evidence for reasoning depth effects
%
% TODO: Game specificity
%   - One coordination game design (invest/arm/attack)
%   - Game parameters systematically varied, but topology is fixed
%
% TODO: Prompt sensitivity
%   - Exact wording matters — report sensitivity analysis from Validation chapter
%   - \cite{lore2024}: framing effects dominate in LLM games
%   - We must distinguish reasoning depth effects from mere framing effects
%
% TODO: Scale
%   - 10-30 agents — large enough for structure, not "society-scale"
%   - Compute constraints limit scale
%
% TODO: Communication conditions
%   - If Phase 2d not completed: acknowledge this as future work
%   - IV3 may be partially tested

% -----------------------------------------------------------------------------
% SECTION 7.5: Implications
% -----------------------------------------------------------------------------
\section{Implications}
% TODO: For computational social science
%   - LLM agents are NOT neutral simulation tools — their reasoning shapes outcomes
%   - Researchers using LLM-ABMs must control for reasoning depth AND information
%   - Memory architecture choice is not "implementation detail" — it changes results
%
% TODO: For AI safety/alignment
%   - Deeper reasoning doesn't mean "better" social outcomes
%   - Non-monotonic effects: more capable AI may produce MORE inequality
%   - Enforcement mechanisms matter: design for commitment, not just capability
%   - \cite{dafoe2020}: our results directly address their open problems
%
% TODO: For mechanism design theory
%   - Empirical evidence that all three Hurwicz inputs matter in LLM systems
%   - Phase transitions may mark qualitative regime changes
%   - Implications for designing multi-agent AI systems
